% Options for packages loaded elsewhere
\PassOptionsToPackage{unicode}{hyperref}
\PassOptionsToPackage{hyphens}{url}
%
\documentclass[
]{article}
\usepackage{amsmath,amssymb}
\usepackage{lmodern}
\usepackage{iftex}
\ifPDFTeX
  \usepackage[T1]{fontenc}
  \usepackage[utf8]{inputenc}
  \usepackage{textcomp} % provide euro and other symbols
\else % if luatex or xetex
  \usepackage{unicode-math}
  \defaultfontfeatures{Scale=MatchLowercase}
  \defaultfontfeatures[\rmfamily]{Ligatures=TeX,Scale=1}
\fi
% Use upquote if available, for straight quotes in verbatim environments
\IfFileExists{upquote.sty}{\usepackage{upquote}}{}
\IfFileExists{microtype.sty}{% use microtype if available
  \usepackage[]{microtype}
  \UseMicrotypeSet[protrusion]{basicmath} % disable protrusion for tt fonts
}{}
\makeatletter
\@ifundefined{KOMAClassName}{% if non-KOMA class
  \IfFileExists{parskip.sty}{%
    \usepackage{parskip}
  }{% else
    \setlength{\parindent}{0pt}
    \setlength{\parskip}{6pt plus 2pt minus 1pt}}
}{% if KOMA class
  \KOMAoptions{parskip=half}}
\makeatother
\usepackage{xcolor}
\usepackage[margin=1in]{geometry}
\usepackage{graphicx}
\makeatletter
\def\maxwidth{\ifdim\Gin@nat@width>\linewidth\linewidth\else\Gin@nat@width\fi}
\def\maxheight{\ifdim\Gin@nat@height>\textheight\textheight\else\Gin@nat@height\fi}
\makeatother
% Scale images if necessary, so that they will not overflow the page
% margins by default, and it is still possible to overwrite the defaults
% using explicit options in \includegraphics[width, height, ...]{}
\setkeys{Gin}{width=\maxwidth,height=\maxheight,keepaspectratio}
% Set default figure placement to htbp
\makeatletter
\def\fps@figure{htbp}
\makeatother
\setlength{\emergencystretch}{3em} % prevent overfull lines
\providecommand{\tightlist}{%
  \setlength{\itemsep}{0pt}\setlength{\parskip}{0pt}}
\setcounter{secnumdepth}{-\maxdimen} % remove section numbering
\ifLuaTeX
  \usepackage{selnolig}  % disable illegal ligatures
\fi
\IfFileExists{bookmark.sty}{\usepackage{bookmark}}{\usepackage{hyperref}}
\IfFileExists{xurl.sty}{\usepackage{xurl}}{} % add URL line breaks if available
\urlstyle{same} % disable monospaced font for URLs
\hypersetup{
  hidelinks,
  pdfcreator={LaTeX via pandoc}}

\author{}
\date{\vspace{-2.5em}}

\begin{document}

\% !TeX program = XeLaTeX

\hypertarget{question}{%
\section{Question}\label{question}}

Los estudiantes de un colegio recolectaron la información
correspondiente a la edad de un grupo de habitantes. La tabla recoge
algunos de los percentiles de las edades de los encuestados.\\

Por ejemplo, el percentil 75 refleja que el 75\% de los encuestados
tiene menos de 62 años. ¿Cuál de las siguientes gráficas representa de
manera correcta los cuartiles de los datos?

\hypertarget{answerlist}{%
\subsection{Answerlist}\label{answerlist}}

\begin{itemize}
\item
  \includegraphics[width=4.5cm,height=\textheight]{image01.png}
\item
  \includegraphics[width=4.5cm,height=\textheight]{image02.png}
\item
  \includegraphics[width=4.5cm,height=\textheight]{image03.png}
\item
  \includegraphics[width=4.5cm,height=\textheight]{image04.png}
\end{itemize}

\hypertarget{solution}{%
\section{Solution}\label{solution}}

El cuartil 1 corresponde al percentil 25 que, de acuerdo con los datos
de la tabla, establece que el 25 \% de los encuestados tiene 26 años de
edad o menos. El cuartil 2 corresponde al percentil 50, que establece
que el 50 \% de los encuestados tiene 43 años de edad o menos.
Finalmente, el cuartil 3 corresponde al percentil 75, que establece que
el 75 \% de los encuestados tiene 62 años de edad o menos. Estos datos
se encuentran representados en la opción:

\begin{tikzpicture}[node distance=5cm, auto,>=latex', thick, scale = 0.5]
  \node{
\begin{tabular}{@{}|l|c|c|c|c|c|c|c|c|c|c|c|c|@{}}
\toprule
\textbf{Cuartil} & 1 & 2 & 3 \\ \midrule
\textbf{Edad}      & 26 & 43 & 62 \\ \bottomrule
\end{tabular}
  };
\end{tikzpicture}

\hypertarget{meta-information}{%
\section{Meta-information}\label{meta-information}}

exname: MMMEdadesTkZ(single-choice) extype: schoice exsolution: 1000
exshuffle: TRUE

\end{document}
